% concatenated from main.tex
\documentclass{article}
\usepackage{graphicx} % Required for inserting images
\usepackage{amsthm}
\usepackage{amsfonts, amsmath, amssymb, amsthm}
\usepackage{cleveref}
\usepackage{url}
\usepackage[utf8x]{inputenc}
\usepackage{cancel}

\author{Quentin Le Hou\'erou \and Ludovic Levy Patey \and Keita Yokoyama}

\title{Conservation of Ramsey's theorem for pairs\\ and well-foundedness}
\date{\today}

\newtheorem{theorem}{Theorem}
\numberwithin{theorem}{section}
\newtheorem{maintheorem}[theorem]{Main Theorem}
\newtheorem{lemma}[theorem]{Lemma}
\newtheorem{question}[theorem]{Question}
\newtheorem{proposition}[theorem]{Proposition}
\newtheorem{statement}[theorem]{Statement}
\newtheorem{remark}[theorem]{Remark}
\newtheorem{definition}[theorem]{Definition}
\newtheorem{corollary}[theorem]{Corollary}

% Allow to repeat theorems
\makeatletter
\newtheorem*{rep@theorem}{\rep@title}
\newcommand{\newreptheorem}[2]{%
\newenvironment{rep#1}[1]{%
 \def\rep@title{#2 \ref{##1}}%
 \begin{rep@theorem}}%
 {\end{rep@theorem}}}
\makeatother

\newreptheorem{maintheorem}{Main Theorem}

% Comments
\usepackage{xcolor}
\newcommand{\ludovic}[1]{\textcolor{blue}{#1}}
\newcommand{\quentin}[1]{\textcolor{green}{#1}}
\newcommand{\keita}[1]{\textcolor{orange}{#1}}

% Reverse maths principles
\newcommand{\RCA}{\mathsf{RCA}}
\newcommand{\WKL}{\mathsf{WKL}}
\newcommand{\SRT}{\mathsf{SRT}}
\newcommand{\EM}{\mathsf{EM}}
\newcommand{\SEM}{\mathsf{SEM}}
\newcommand{\ADS}{\mathsf{ADS}}
\newcommand{\SCAC}{\mathsf{SCAC}}
\newcommand{\SADS}{\mathsf{SADS}}
\newcommand{\COH}{\mathsf{COH}}
\newcommand{\ACA}{\mathsf{ACA}}

% Forcing question
\newcommand{\qvdash}{\operatorname{{?}{\vdash}}}
\newcommand{\nqvdash}{\operatorname{{?}{\nvdash}}}

% Induction principles
\newcommand{\BSig}{\mathsf{B}\Sigma^0}
\newcommand{\ISig}{\mathsf{I}\Sigma^0}
\newcommand{\LSig}{\mathsf{L}\Sigma^0}
\newcommand{\BPi}{\mathsf{B}\Pi^0}
\newcommand{\IPi}{\mathsf{I}\Pi^0}
\newcommand{\LPi}{\mathsf{L}\Pi^0}
\newcommand{\IDel}{\mathsf{I}\Delta^0}
\newcommand{\BCDe}{\mathsf{BC}\Delta^0}
\newcommand{\WF}{\mathsf{WF}}
\newcommand{\RT}{\mathsf{RT}}
\newcommand{\BME}{\mathsf{BME}}
\newcommand{\TT}{\mathsf{TT}}

\renewcommand{\O}{\mathcal{O}}
\renewcommand{\L}{\mathcal{L}}
\newcommand{\C}{\mathcal{C}}
\newcommand{\R}{\mathcal{R}}
\renewcommand{\S}{\mathcal{S}}
\newcommand{\T}{\mathcal{T}}
\renewcommand{\P}{\mathcal{P}}
\newcommand{\M}{\mathcal{M}}
\newcommand{\Nc}{\mathcal{N}}
\renewcommand{\P}{\mathcal{P}}
\newcommand{\Q}{\mathcal{Q}}
\newcommand{\NN}{\mathbb{N}}
\newcommand{\PP}{\mathbb{P}}
\newcommand{\uh}[0]{{\upharpoonright}}


\newcommand{\card}{\operatorname{card}}



% Natural operations on ordinals
\usepackage{mathabx}
\newcommand{\nplus}{\dotplus}
\newcommand{\ntimes}{\dottimes}
\newcommand{\nminus}{\dotdiv}
\newcommand{\nsum}{\dot{\sum}}

% \newcommand{\nplus}[0]{\#}
% \newcommand{\ntimes}[0]{\smashtimes}
% % \newcommand{\nsum}[0]{\mathlarge{\mathlarger{\#}}}
% \DeclareMathOperator*{\nsum}{\text{\LARGE$ \# $}}

%typographically correct double quotes

\def\qt#1{``#1''}%

\title{Shore blocking for well-foundedness: an infinite descending sequence below the natural product of the ordinal bound with k}
\author{Quentin Le Houérou, Ludovic Levy Patey, Keita Yokoyama}
\date{arxiv:2402.11616v3; CC BY 4.0}
\begin{document}
\maketitle
\setcounter{section}{2}
\setcounter{theorem}{2}
\noindent\textit{Editorial note.} The counted statement is the lemma labelled lem:shore-blocking-wf at source lines 426-431, and the proof environment at source lines 433-450 is its own proof by adjacency: it is the only proof environment attached to that label and the lemma is not restated anywhere else in the article.\par\medskip
\noindent\textit{Editorial note.} The retained context is the block of the same subsection that the counted lemma depends on and belongs with: the definition of the natural sum and natural product at source lines 392-398, the strict monotonicity proposition at source lines 402-404 with its proof at 406-408, and the companion lemma lem:wf-product at source lines 411-413 with its proof at 415-423. With the section and theorem counters set to two and two these print as Definition 2.3, Proposition 2.4, Lemma 2.5 and Lemma 2.6, exactly the numbers they carry in the article.\par\medskip
\noindent\textit{Editorial note.} The article's own preamble is kept, including cleveref, which the article uses for its bracketed label syntax, and mathabx, which supplies the natural sum and natural product symbols. One block only is left out: the six lines at source lines 93 to 99 that declare the LS2 font encoding, substitute the stix family for it and define the single binary symbol built on it, because the production engine's TeX installation carries no stix font definition file and the declaration therefore aborts the run before the first page. That symbol is used nowhere in the delivered excerpt and appears only inside a commented-out macro in the retained preamble, so no symbol of the counted statement or of its proof is lost; no package is added, no macro is invented and no author text is altered.\par\medskip
\noindent\textit{Editorial note.} No retained line contains a citation command, a figure, an included graphics file or a reference to a label defined outside the retained spans, so the package declares no bibliography, no asset and no external reference, and no unresolved reference or citation can arise.\par\medskip
\noindent\textit{Editorial note.} Both retained lemmas and the retained proposition are delivered line for line exactly as printed, including the trailing sentence of the counted statement that an index of the blocking functional can be found computably in k.\par\medskip
\noindent\textit{Editorial scope.} One original result: the lemma labelled lem:shore-blocking-wf (source lines 426-431) with its complete adjacent author proof (source lines 433-450), together with the author's definition of the natural sum and natural product (source lines 392-398), the strict monotonicity proposition (source lines 402-408) and the companion lemma lem:wf-product (source lines 411-423) of the same subsection. The forcing lemma lem:coh-wf-forcing-wf at source lines 522-524 with its proof at 525-543, the conservation theorems of the article and the later Sections 3 to 5 are not counted here and are not used as premises.\par\medskip
\noindent\textit{Allowed standard prerequisites (editorial).}\begin{itemize}
\item Second-order arithmetic, countable topped models and the principle WF(alpha). RCA\_0 is the usual base theory of reverse mathematics; a countable model M = (M,S) is topped by a set Y in S when every set of the model is computable from Y; WF(alpha) states that the ordinal alpha is well founded, and an M-infinite decreasing sequence of ordinals below a bound is a sequence indexed by the model's own natural numbers.. Fixes the meaning of the hypothesis of the counted lemma and of the conclusion that a Turing functional is an M-infinite decreasing sequence below the natural product bound.
\item Turing functionals, indices and their computable manipulation. Turing functionals with oracle X are written as Phi-sub-e to the X, an index of a functional can be found computably from its parameters, and functionals may be twisted so that the value at a given input depends only on an initial segment of the oracle.. The proof begins by twisting the functionals so that convergence forces the three normalisation properties, and the statement itself promises that an index of the blocking functional can be found computably in k.
\item Natural (Hessenberg) sum and product below epsilon-zero. For ordinals below epsilon-zero written in Cantor normal form with a common family of exponents, the natural sum adds the coefficients exponent by exponent, the natural product alpha times k is the natural sum of k copies of alpha, and the natural sum is strictly increasing in each argument.. The blocking functional outputs the natural sum of the k last values, and the strict decrease of that sum is what makes the output an M-infinite descending sequence below the natural product bound.
\end{itemize}\medskip
\begin{definition}[Natural sum and product]
Let $\alpha$ and $\beta$ be two ordinals less than $\epsilon_0$. Let $\alpha = \omega^{\gamma_1}n_1 + \dots + \omega^{\gamma_k}n_k$ and $\beta = \omega^{\gamma_1}m_1 + \dots + \omega^{\gamma_k}m_k$ (We allow the $n_i$ and $m_i$ to be equal to $0$ in order to write $\alpha$ and $\beta$ using the same exponents $\gamma_i$). We define the natural sum of $\alpha$ and $\beta$ (denoted by $\alpha \nplus \beta$) to be the following ordinal (It does not depend on the choice of the family $(\gamma_i)$) : 

$$\omega^{\gamma_1}(n_1 + m_1) + \dots + \omega^{\gamma_k}(n_k + m_k)$$

We also define $\alpha \ntimes k$ to be equal to be the natural sum of $\alpha$ with itself $k$ times. %and $\alpha \ntimes \omega = \omega^{\gamma_1+1}n_1 + \dots + \omega^{\gamma_k+1}n_k$ (whereas $\alpha \times \omega = \omega^{\gamma_1 + 1}n_1$)
\end{definition}

\begin{proposition}
For all $\beta < \epsilon_0$, and $\alpha' < \alpha < \epsilon_0$ : $\alpha' \nplus \beta < \alpha \nplus \beta$ and $\beta \nplus \alpha' < \beta \nplus \alpha$.
\end{proposition}

\begin{proof}
Clear from the definition of $\nplus$.
\end{proof}

\begin{lemma}\label[lemma]{lem:wf-product}
Let~$\M = (M,S) \models \RCA_0 + \WF(\alpha)$ for $\alpha$ an ordinal less than $\epsilon_0$. Then $\M \models \WF(\alpha \ntimes k)$ for every $k \in M$.
\end{lemma}

\begin{proof}
Since $\alpha \ntimes k \leq \alpha \times \omega$ for every $k \in M$ and $\RCA_0$ proves that the product of two well-orders is a well-order, hence $\RCA_0 \vdash \WF(\alpha) \rightarrow \WF(\alpha \times \omega)$, we get the desired result.

% Consider an M-infinite decreasing sequence of ordinals $(\alpha_s)_{s \in M}$ all less than $\alpha \ntimes k$. Each $\alpha_s$ can be written as $\alpha_{s,0} \# \dots \# \alpha_{s,k-1}$ where we choose $\alpha_{s,0}$ to be the biggest ordinal less than or equal to $\alpha$ satisfying $\alpha_{s,0} \leq \alpha_s$, we then choose $\alpha_{s,1}$ to be the biggest ordinal less than or equal to $\alpha$  satisfying $\alpha_{s,0} \# \alpha_{s,1} \leq \alpha_s$...

% Since $(\alpha_s)$ is decreasing, the sequences $(\alpha_{s,i})$ are non-decreasing for every $i < k$, and for every $s$, there is at least one $i < k$ such that $\alpha_{s,i} > \alpha_{s+1,i}$.

% But, by $\RT^1_k$, which follows from $\BSig_2$, there exists some~$i < k$ and an $M$-infinite set~$H \in \M$ such that $\alpha_{s,i} > \alpha_{s+1,i}$ for all $s \in H$. Then $\{\alpha_{s,i} : s \in H \}$ is an $M$-infinite decreasing sequence in~$\alpha$, contradicting our assumption.
\end{proof}

\begin{lemma}\label[lemma]{lem:shore-blocking-wf}
Consider a model $\M = (M,S)$ and $\alpha \in M$ an ordinal less than $\epsilon_0$.
For every~$k \in M$, there is a Turing functional $\Gamma_k$ such that, letting~$\beta \in \alpha$ be the largest ordinal less than $\alpha$ of index at most~$k$, for every~$X \in 2^M$ such that $\M[X] \models \RCA_0$, if there is some~$e < k$ such that $\Phi^X_e$ is an $M$-infinite decreasing sequence of elements smaller than~$\beta$, then $\Gamma^X_k$ is an $M$-infinite decreasing sequence of elements smaller than~$\beta \ntimes k$.

Moreover, an index of~$\Gamma_k$ can be found computably in~$k$.
\end{lemma}

\begin{proof}
By twisting the Turing functionals, we can assume that for every~$e, n \in M$, if $\Phi^\sigma_e(n)\downarrow$, then
\begin{itemize}
    \item[(1)] $n < |\sigma|$ ;
    \item[(2)] $\Phi^\sigma_e(m)\downarrow$ for every~$m < n$ ;
    \item[(3)] $\Phi^\sigma_e(0), \Phi^\sigma_e(1), \dots, \Phi^\sigma_e(n)$ is a strictly decreasing sequence of elements smaller than~$\beta$.
\end{itemize}

Given $\sigma \in 2^{<M}$ and $e < k$, let $\alpha(\sigma, e) = \Phi^{\sigma}_e(s)$ for the largest $s < |\sigma|$ such that $\Phi^{\sigma}_e(s)\downarrow$. If there is no such~$s$, then $\alpha(\sigma, e) = \beta$. Note that if $\sigma' \succeq \sigma$, then $\alpha(\sigma',e) \leq \alpha(\sigma,e)$.

Let~$\sigma_{-1} = \epsilon$.
Let~$\Gamma^X_k$ be the Turing functional which, on input~$n$, searches for some~$x > |\sigma_{n-1}|$ and some~$\sigma_n \prec X$ such that $\Phi_e^{\sigma_n}(x)\downarrow$ for some~$e < k$. If found, it outputs $\alpha(\sigma, 0) \nplus \dots \nplus \alpha(\sigma, k-1)$. Note that if~$\Gamma^X_k(n)\downarrow$, then by (3), $\Gamma^X(n)$ is an ordinal smaller than $\beta \ntimes k$.

Suppose that $X$ is such that $\M[X] \models \RCA_0$ and there is an $e < k$ is such that $\Phi_e^X$ is total.
Since~$\M[X] \models \ISig_1$, $\Gamma^X$ is total.

Moreover, since~$x > |\sigma_{n-1}|$, letting~$e < k$ be such that $\Phi_e^{\sigma_n}(x)\downarrow$, by (1) we have $\Phi_e^{\sigma_{n-1}}(x)\uparrow$, so by (2) and (3), $\alpha(\sigma_{n+1},e) < \alpha(\sigma_n, e)$, hence $\Gamma^X_k(n+1) < \Gamma^X_k(n)$. It follows that $\Gamma^X_k$ is an $M$-infinite decreasing sequence of ordinals smaller than~$\beta \ntimes k$.
\end{proof}

\end{document}
